\head{第四卷勘误清单三：与第0章§§16--19已译部分有关的选译}
\source{EGA IV第四部分，勘误清单三，印刷第345、347页；上传PDF第344、346页}
\NewNote{本节只收入与已译范围直接相关的七项作者勘误，保留其原分类和原编号。它不是整份勘误清单三的全译；未列出的条目不因此被删除或判定无关。下文的“页”均指作者所列原出版分册的印刷页码。}

\egasection{A）排印错误：三项选译}
\hypertarget{ErrIV:A16222}{}\noindent\textbf{（$0_{\mathrm{IV}}$，16.2.2.2）。}
第26页第13行，将 $M'\cap\mathfrak q^nM'$ 改为 $M'\cap\mathfrak q^nM$。

\medskip
\hypertarget{ErrIV:A1737}{}\noindent\textbf{（$0_{\mathrm{IV}}$，17.3.7）。}
第50页倒数第16行，将（17.1.10）改为（17.3.3）。

\medskip
\hypertarget{ErrIV:A1813}{}\noindent\textbf{（$0_{\mathrm{IV}}$，18.1.3）。}
在图（18.1.3.1）中，$p_2$ 应标在箭头 $G\to F$ 的旁边。

\egasection{B）正文修订：四项选译}
\hypertarget{ErrIV:5}{}
\noindent\textbf{（Err$_{\mathrm{IV}}$，5）。}
在（$0_{\mathrm{IV}}$，16.1.4）中，第23页倒数第8行，将“说 $A$ 是悬链的”改为：“说一个整环 $A$ 是悬链的”。

\medskip
\hypertarget{ErrIV:6}{}
\noindent\textbf{（Err$_{\mathrm{IV}}$，6）。}
在（$0_{\mathrm{IV}}$，16.1.9）中，可删去 $A$ 为诺特环的假设。事实上，对于每个 $u\in B_M$，$A[u]$ 是 $C$ 的一个子 $A$-代数，故 $u$ 整于 $A$（Bourbaki，\textit{交换代数}，第五章，§1，第1目，命题1的推论1），因而也整于 $A_M$；仍由（16.1.5）得出结论。

\medskip
\hypertarget{ErrIV:7}{}
\noindent\textbf{（Err$_{\mathrm{IV}}$，7）。}
在（$0_{\mathrm{IV}}$，16.5.12）中，应注意：关系（16.5.11.1）等价于将 $A$ 换成各商环 $A/\mathfrak r$ 后得到的相同关系，其中 $\mathfrak r$ 遍历 $A$ 的极小素理想。因此，在（16.5.12）的证明中可限于 $A$ 为整环的情形，并应用（16.1.4.2）。


\medskip
\hypertarget{ErrIV:8}{}
\noindent\textbf{（Err$_{\mathrm{IV}}$，8）。}
删去（$0_{\mathrm{IV}}$，19.3.12），因为其证明错误（集合 $W_\lambda,U_\lambda,V_\lambda$ 不是线性流形）。
